\documentclass[a4paper]{article}
\usepackage{amsfonts}
\usepackage{amsmath}
\usepackage{amssymb}
\usepackage{graphicx}  
\usepackage{cancel}

\newcommand{\degree}{^{\circ}}
\newcommand{\cis}{\operatorname{cis}}
\newcommand{\Bgsin}{\operatorname{Bgsin}}
\newcommand{\Bgcos}{\operatorname{Bgcos}}
\newcommand{\Bgtan}{\operatorname{Bgtan}}
\newcommand{\limit}{\lim\limits}
\newcommand{\summation}{\displaystyle\sum}
\newcommand{\dom}{\operatorname{dom}}
\newcommand{\Int}{\displaystyle\int}

\begin{document}
  
\noindent \large Auteur: Rune Suy \\
\noindent \large Studierichting: Informatica\\
\noindent \large Oefeningenreeks X nr. Y\\

\medskip

\normalsize

$f(x,y) = 2x^3 + y^3 - 3x^2 - 12x - 3y$\\

$\dfrac{\partial f}{\partial x} = 6x^2 - 6x - 12$\\

$\dfrac{\partial f}{\partial y} = 3y^2 - 3$\\

$\left\{
    \begin{array}{l}
        6x^2 - 6x - 12 = 0\\
        3y^2 - 3 = 0
    \end{array}
\right.$\\

$\Rightarrow y=1$ of $y=-1$\\

$D=\sqrt{18}$\\

$x_1 = -1$\\

$x_2 = 2$\\

$H(x,y) = \det \left(\begin{array}{cc}
    12x-6 & 0\\
    0 & 6y
\end{array}\right) = 72xy-36y$\\

$H(-1,1) = -72-36 < 0$ dus $(-1,1)$ is een zadelpunt\\

$H(-1,-1) = 72+36 > 0$ en $\dfrac{\partial^2 f}{\partial x^2}, \dfrac{\partial^2 f}{\partial y^2} > 0$ dus $(-1,-1)$ is een maximum\\

$H(2,-1) = -144 - 36 < 0$ dus $(2,-1)$ is een zadelpunt\\

$H(2,1) = 72 - 36 > 0$ en $\dfrac{\partial^2 f}{\partial x^2}, \dfrac{\partial^2 f}{\partial y^2} < 0$ dus $(2,1)$ is een minimum\\

\begin{thebibliography}{999}
%\bibitem{a} Referentie
\end{thebibliography}
\end{document} 
